\documentclass[11pt,a4paper]{article}
\usepackage[T1]{fontenc}
\usepackage[utf8]{inputenc}
\usepackage[french]{babel}
\usepackage[margin=2.2cm]{geometry}
\usepackage{booktabs}
\usepackage{array}
\usepackage{tabularx}
\usepackage{csquotes}
\usepackage{xcolor}
\usepackage{colortbl}
\usepackage{enumitem}
\usepackage{hyperref}

\definecolor{flag}{RGB}{170,40,40}
\definecolor{ok}{RGB}{36,130,70}
\definecolor{mix}{RGB}{180,120,20}
\definecolor{tableRowAlt}{RGB}{238,242,248}
\newcolumntype{Y}{>{\raggedright\arraybackslash}X}
\newcommand{\hum}[1]{\textcolor{ok}{#1}}
\newcommand{\iaa}[1]{\textcolor{flag}{#1}}
\newcommand{\mixx}[1]{\textcolor{mix}{#1}}

\title{Audit stylométrique IA / Humain --- \textbf{méthode v2 (contrastive)}\\
\large Fichier \texttt{main.tex} --- approfondissement}
\author{Auditeur automatique (Claude Code)}
\date{\today}

\begin{document}
\maketitle

\section{Ce que corrige cette version}
La première passe rendait un global $\approx$ 50\,\%, peu exploitable. \textbf{Ce n'était
pas une erreur de mesure mais un artefact de méthode :} moyenner des signaux faibles en
valeur \emph{absolue} ramène mécaniquement tout vers le centre (régression vers la
moyenne). Deux changements corrigent cela :
\begin{enumerate}[]
  \item \textbf{Audit contrastif.} Le document contient des sections \emph{prouvées
  humaines} (présentation Geismar, descriptif du poste : 1\textsuperscript{re} personne,
  faits vérifiables, zéro tournure générique). Elles servent d'\textbf{empreinte de
  référence de l'auteur}. On mesure la \emph{distance} de chaque section à cette empreinte,
  au lieu d'un seuil absolu universel.
  \item \textbf{Features réellement discriminantes} (et non la moyenne des longueurs) :
  \emph{forme} de la distribution des longueurs de phrase (l'IA évite les phrases
  courtes et varie peu), \emph{diversité des ouvreurs}, \emph{répétition de n-grammes} ;
  plus deux signaux quasi-binaires : densité de \emph{coquilles réelles} (marqueur humain)
  et densité \emph{1\textsuperscript{re} personne + concret} (chiffres, identifiants de code).
\end{enumerate}
Résultat : les scores s'étalent désormais de \textbf{12 à 79} (et non plus 38--62), et
deux indicateurs indépendants (rubrique + distance à l'empreinte) concordent.

\section{Limite assumée}
La littérature identifie la \textbf{perplexité} (régularité probabiliste sous un modèle de
langage) comme le signal le plus discriminant. Elle exige un modèle de langage : je ne la
calcule \textbf{pas} et ne la simule pas (cela serait inventé). La méthode ci-dessous est
\emph{stylométrique} et \emph{contrastive} ; elle reste heuristique, non forensique.

\section{Diagnostic global révisé}
Le $\approx$ 49\,\% global n'est pas \enquote{aléatoire} : c'est la signature d'un document
\textbf{Mixte par construction}, où la plupart des sections combinent un \emph{fond humain
réel} et une \emph{forme rédactionnelle IA}. La décomposition par volume de prose le montre :

\begin{center}
\begin{tabularx}{\linewidth}{Y r l}
\toprule
Bande de score (sur 9\,142 mots de prose) & Part & Interprétation \\
\midrule
\hum{$\leq$ 35 --- Humain probable} & $\approx$ 16\,\% & Voix perso, concret, coquilles \\
36--44 --- Mixte équilibré & $\approx$ 9\,\% & Indécis \\
\mixx{45--59 --- Mixte à dominante IA} & $\approx$ 73\,\% & Fond réel, \emph{forme} IA \\
\iaa{$\geq$ 60 --- IA probable} & $\approx$ 2\,\% & Synthèses de clôture, généralités \\
\bottomrule
\end{tabularx}
\end{center}

\noindent \textbf{Conclusion :} environ \textbf{trois quarts de la prose} sont du
\enquote{contenu réel rédigé dans une forme IA}, avec un noyau personnel net ($\approx$
16\,\%) et de rares passages purement génériques. Ce n'est pas un coup de dé : c'est un
document rédigé \textbf{avec} une IA sur une matière authentique.

\section{Empreinte humaine de référence (auteur)}
Calculée sur \enquote{Geismar} + \enquote{Le poste} :
burstiness (CV) $\approx$ 0,44 ; phrases courtes $\approx$ 35\,\% ; diversité des ouvreurs
$\approx$ 0,90 ; \enquote{tells} $\approx$ 0 ; densité concret $\approx$ 85/1000.
Toute section qui s'éloigne fortement de ce profil (CV bas, peu de phrases courtes,
ouvreurs répétitifs, concret faible) est suspecte.

\section{Résultats par section (score IA 5--95, décisif)}
\begin{center}\small
\begin{tabularx}{\linewidth}{Y r c c Y}
\toprule
Section & Mots & Score & Verdict & Signal moteur \\
\midrule
\rowcolor{tableRowAlt}\multicolumn{5}{l}{\textit{Pages liminaires}}\\
Remerciements & 112 & 39 & Mixte & tells 17,9 vs coquilles+\enquote{je}$\times$9 \\
Introduction & 257 & 31 & \hum{Humain} & rythme varié (CV 0,63) ; \emph{mais 1\textsuperscript{re} phrase IA} \\
Entreprise : Geismar & 108 & 28 & \hum{Humain} & faits concrets, ancre de référence \\
Le poste (alternance) & 171 & 12 & \hum{Humain} & \enquote{je}, concret, ouvreurs 1,0 \\
Synthèse des blocs & 70 & 36 & Mixte & court (faible confiance) \\
\midrule
\rowcolor{tableRowAlt}\multicolumn{5}{l}{\textit{Bloc 1 --- Challenge Datalake}}\\
Présentation projets & 1515 & \mixx{59} & \mixx{Mixte (IA)} & ouvreurs \iaa{0,44}, peu de phrases courtes \\
Justification compét. & 1197 & 48 & Mixte & ouvreurs \iaa{0,38}, clôtures répétées \\
Bilan bloc 1 & 106 & 48 & Mixte & court (faible confiance) \\
\midrule
\rowcolor{tableRowAlt}\multicolumn{5}{l}{\textit{Bloc 2 --- Spark \& DevOps}}\\
Présentation projets & 2295 & \mixx{56} & \mixx{Mixte (IA)} & ouvreurs \iaa{0,25} (pire) ; mais concret+coquilles \\
\midrule
\rowcolor{tableRowAlt}\multicolumn{5}{l}{\textit{Bloc 3 --- Cloud AWS}}\\
Contexte et enjeux & 157 & 32 & \hum{Humain} & CV 0,50, phrases courtes 25\,\% \\
Architecture cloud & 101 & 34 & \hum{Humain} & concret 99, varié \\
Stockage S3 & 108 & 40 & Mixte & — \\
DynamoDB \& GSI & 226 & 30 & \hum{Humain} & concret 120, chiffres, coquille \\
Lambda \& API Gateway & 172 & 20 & \hum{Humain} & CV 0,51, concret 116 \\
Sécurité \& IAM & 155 & 40 & Mixte & — \\
Optimisation coûts & 119 & 39 & Mixte & très proche de l'empreinte humaine \\
Synthèse Bloc 3 & 44 & \iaa{61} & \iaa{IA prob.} & CV 0,09 ; \enquote{couvre l'intégralité} \\
\midrule
\rowcolor{tableRowAlt}\multicolumn{5}{l}{\textit{Bloc 4 --- Machine Learning}}\\
Compétences visées & 88 & 54 & \mixx{Mixte (IA)} & CV 0,11 (énumération mécanique) \\
Contexte et enjeux & 113 & 43 & Mixte & — \\
Architecture / cycle ML & 88 & \iaa{62} & \iaa{IA prob.} & CV \iaa{0,12} ; \enquote{il est impératif} \\
Analyse exploratoire (EDA) & 305 & 25 & \hum{Humain} & 24 chiffres, CV 0,57, courtes 35\,\% \\
Préparation des données & 196 & 58 & \mixx{Mixte (IA)} & tells 5,1 ; tirets 10 \\
Entraînement comparatif & 245 & 56 & \mixx{Mixte (IA)} & ouvreurs 0,64 ; tells 4 \\
Résultats et évaluation & 494 & 46 & Mixte & rép. n-gr. 47 \emph{mais} 89 chiffres \\
Mise en service & 251 & 53 & Mixte & tells 4 \\
Robustesse \& limites & 190 & 54 & \mixx{Mixte (IA)} & CV 0,23 ; \enquote{d'une part\,\dots} \\
Synthèse Bloc 4 & 52 & \iaa{79} & \iaa{IA prob.} & CV \iaa{0,04} ; rép. n-gr. \iaa{58} \\
\bottomrule
\end{tabularx}
\end{center}
\footnotesize Confiance : élevée pour les sections $>$ 250 mots ; faible $<$ 120 mots
(synthèses, chapeaux). \enquote{Verdict} dérivé du score : \hum{$\leq$35 Humain},
36--44 Mixte, \mixx{45--59 Mixte (IA)}, \iaa{$\geq$60 IA}.
\normalsize

\section{Classement actionnable}
\textbf{À personnaliser en priorité (les plus \enquote{IA}) :}
Synthèse Bloc 4 (79) ; Architecture / cycle ML (62) ; Synthèse Bloc 3 (61) ;
Bloc 1 présentation (59) ; Préparation données (58) ; Bloc 2 présentation \& Entraînement (56).

\smallskip
\textbf{À conserver tel quel (les plus humains) :}
Le poste (12) ; Lambda \& API (20) ; EDA (25) ; Geismar (28) ; DynamoDB (30).

\section{Signaux les plus décisifs (mesurés)}
\begin{center}
\begin{tabularx}{\linewidth}{Y c c Y}
\toprule
Signal & Sections humaines & Sections IA & Pouvoir séparateur \\
\midrule
Diversité des ouvreurs & 0,80--1,00 & 0,25--0,50 & \textbf{fort} \\
\% phrases courtes ($<$11 mots) & 25--40\,\% & 0--8\,\% & \textbf{fort} \\
Burstiness (CV) & 0,40--0,63 & 0,04--0,23 & fort \\
Densité concret (chiffres/code) & 85--200 & 12--40 & moyen \\
Coquilles réelles & présentes & absentes & moyen (binaire) \\
\enquote{Tells} & $\approx$ 0 & variable & faible (sauf liminaires) \\
\bottomrule
\end{tabularx}
\end{center}
\noindent Contrairement aux idées reçues, le \emph{lexique} (\enquote{tells}) discrimine
mal ici : le signal fort est \textbf{structurel} (rythme des phrases, ouvreurs). C'est
cohérent avec la littérature : la longueur moyenne des phrases et les bigrammes de
mots-outils sont les marqueurs les plus discriminants, l'IA produisant des phrases plus
uniformes (peu de phrases courtes) que l'humain.

\section{Zones non rédigées (inchangé, bloquant)}
Bloc 5 entier, Conclusion, Bloc 2 \enquote{Analyse attendue}, page de titre
(\enquote{A completer}, \texttt{\textbackslash author\{Tonio\}}) : texte de gabarit, exclu
du scoring. À traiter avant dépôt. Détails et trame dans \texttt{propositions\_reecriture.tex}.

\section{Verdict}
\textbf{Document Mixte}, fond authentique et personnel, forme rédactionnelle largement
assistée par IA (signature structurelle nette sur $\approx$ 73\,\% de la prose). Noyau
clairement humain de $\approx$ 16\,\% (alternance, EDA, accès DynamoDB/Lambda). Pour
abaisser l'empreinte IA, agir sur le \textbf{rythme} (réintroduire des phrases courtes,
diversifier les ouvreurs) plutôt que sur le vocabulaire, et personnaliser les synthèses
de clôture. Estimation indicative, à valider par jugement humain.

\section*{Sources (méthode)}
\small
\begin{itemize}
  \item Pangram Labs --- \emph{Why Perplexity and Burstiness Fail to Detect AI}.
  \url{https://www.pangram.com/blog/why-perplexity-and-burstiness-fail-to-detect-ai}
  \item \emph{Feature-Based Detection of AI-Generated Text: Stylometric and Perplexity
  Markers} (ResearchGate, 2025).
  \item \emph{Stylometric comparisons of human versus AI-generated creative writing},
  Humanities and Social Sciences Communications (Nature), 2025.
  \url{https://www.nature.com/articles/s41599-025-05986-3}
  \item \emph{A Lightweight Approach to Detection of AI-Generated Text}, arXiv 2511.21744.
\end{itemize}
\normalsize
\end{document}
